\subsubsection{Subproceso 2: Selección y Registro de Cursos}

\textbf{Fuentes documentales:}
\begin{itemize}
    \item \textbf{Resolución Rectoral N.\textsuperscript{o} RR-0570-2022 (UNI):} Reglamento de Matrícula que define la tabla de créditos máximos según el promedio ponderado del estudiante.
    \item \textbf{TUPA de la UNI --- Procedimiento PA5400954D (RR-3698-2024):} Documento oficial que establece que la Matrícula Regular es de aprobación automática e indica sus requisitos.
    \item \textbf{Estatuto de la UNI y Reglamento de Matrícula vigente:} Normas de la universidad que regulan las condiciones para inscribir cursos, la prioridad de asignaturas desaprobadas y el riesgo académico.
    \item \textbf{Ley Universitaria N.\textsuperscript{o} 30220:} Ley general que regula los estudios universitarios en el Perú.
\end{itemize}

\textbf{Fuentes tecnológicas:}
\begin{itemize}
    \item \textbf{Intranet DIRCE (\url{https://dirce.uni.edu.pe/}):} Portal de matrícula revisado directamente para observar los pasos de selección de cursos, las alertas en pantalla y las validaciones automáticas.
    \item \textbf{Sistema SIGA-ORCE:} Sistema central de gestión académica cuyo funcionamiento se analizó a partir de las respuestas del portal web y de los reglamentos de la universidad.
\end{itemize}

\textbf{Fuentes humanas:}
\begin{itemize}
    \item \textbf{Estudiantes regulares de la UNI:} Alumnos consultados para comparar lo que indican las normas con su experiencia real usando el portal durante la matrícula, sobre todo respecto a la lentitud del sistema y el rápido agotamiento de vacantes.
\end{itemize}
